% AUTO-GENERATED by scripts/generate_bilingual_glossary.py. DO NOT EDIT.
\subsection{全量中英对照：149 项术语速查}
\phantomsection\label{app:bilingual-glossary}

这张表按 English 名称首字母分为 A--D 至 U--Z 六组；
中英文并排列出基础与现代概念。
先看一句话含义，再按最右侧节号／页码查找正文中的解释、公式或算例；例如 20.3／59 依次表示节号与页码。
中文查找可跳到\hyperref[app:chinese-index]{中文拼音索引}。

\begingroup\small
\setlength{\tabcolsep}{3pt}
\renewcommand{\arraystretch}{1.18}
\tagpdfsetup{table/header-rows={1}}
\begin{longtable}{@{}>{\raggedright\arraybackslash}p{38mm}>{\raggedright\arraybackslash}p{28mm}>{\raggedright\arraybackslash}p{62mm}>{\centering\arraybackslash}p{17mm}@{}}
\toprule
English & 中文 & 一句话定义 & 节号／页码\\\midrule
\endfirsthead
\toprule
English & 中文 & 一句话定义 & 节号／页码\\\midrule
\endhead
\bottomrule
\endfoot
\addlinespace[4pt]
\multicolumn{4}{@{}l}{\sffamily\bfseries\color{blue} English A--D / 英文首字母索引}\\*
Activation function & 激活函数 & 应用在线性层输出上的函数，扩展模型表达能力。 & \hyperref[intro:07-functions-relu]{\ref*{intro:07-functions-relu}/\pageref*{intro:07-functions-relu}}\\
Adam & Adam 优化器 & 结合动量和自适应尺度估计的常用优化器。 & \hyperref[intro:18c-training-engineering]{\ref*{intro:18c-training-engineering}/\pageref*{intro:18c-training-engineering}}\\
AdamW & AdamW 优化器 & 保留 Adam 的自适应更新并独立施加权重衰减。 & \hyperref[concept:adamw]{\ref*{concept:adamw}/\pageref*{concept:adamw}}\\
Agent & 智能体 & 围绕目标选择工具、执行动作并利用结果继续决策的系统。 & \hyperref[concept:agent]{\ref*{concept:agent}/\pageref*{concept:agent}}\\
Application memory & 应用记忆 & 应用保存并在需要时恢复的信息，不等于模型权重或 KV Cache。 & \hyperref[concept:app-memory]{\ref*{concept:app-memory}/\pageref*{concept:app-memory}}\\
Associative memory interference & 联想记忆串扰 & 多个不正交 Key 的读取贡献叠加，污染目标 Value。 & \hyperref[concept:memory-interference]{\ref*{concept:memory-interference}/\pageref*{concept:memory-interference}}\\
Autoregressive generation & 自回归 & 每次以前缀为条件预测下一 Token，并把选定 Token 接回上下文。 & \hyperref[concept:autoregressive]{\ref*{concept:autoregressive}/\pageref*{concept:autoregressive}}\\
Backpropagation & 反向传播 & 利用链式法则高效计算计算图中全部参数的梯度。 & \hyperref[intro:17-backpropagation]{\ref*{intro:17-backpropagation}/\pageref*{intro:17-backpropagation}}\\
Batch size & 批量大小 & 一个批次包含多少训练样本。 & \hyperref[intro:18c-training-engineering]{\ref*{intro:18c-training-engineering}/\pageref*{intro:18c-training-engineering}}\\
Bias & 偏置 & 线性变换后添加的常数或向量，使模型不必经过原点。 & \hyperref[intro:10-mlp]{\ref*{intro:10-mlp}/\pageref*{intro:10-mlp}}\\
BOS/EOS & 开始与结束标记 & 词表中的开始或结束标记，使用方式和结束 ID 由模型配置规定。 & \hyperref[concept:bos-eos]{\ref*{concept:bos-eos}/\pageref*{concept:bos-eos}}\\
Byte pair encoding (BPE) & 字节对编码 & 在语料中反复合并高频相邻 Token 对，建立子词单元。 & \hyperref[concept:bpe]{\ref*{concept:bpe}/\pageref*{concept:bpe}}\\
Causal mask & 因果掩码 & 防止生成位置注意到未来 Token 的约束。 & \hyperref[concept:causal-mask]{\ref*{concept:causal-mask}/\pageref*{concept:causal-mask}}\\
Chain-of-thought (CoT) prompting & 思维链提示 & 可能要求模型输出中间步骤的提示形式，不等同于直接观察内部计算。 & \hyperref[concept:cot-prompt]{\ref*{concept:cot-prompt}/\pageref*{concept:cot-prompt}}\\
Chat template & 聊天模板 & 把消息角色与内容序列化为模型配套的 Token 输入格式。 & \hyperref[concept:chat-template]{\ref*{concept:chat-template}/\pageref*{concept:chat-template}}\\
Context & 上下文 & 一次推理可读取的对话、检索片段和工具结果等上下文内容。 & \hyperref[concept:context]{\ref*{concept:context}/\pageref*{concept:context}}\\
Context window & 上下文窗口 & 单次处理所允许使用的 Token 长度范围。 & \hyperref[intro:23-llm-systems]{\ref*{intro:23-llm-systems}/\pageref*{intro:23-llm-systems}}\\
Convolution kernel & 卷积核 & 在图片局部区域滑动的可训练权重小矩阵。 & \hyperref[intro:19-cnn]{\ref*{intro:19-cnn}/\pageref*{intro:19-cnn}}\\
Convolutional neural network & 卷积神经网络 & 卷积神经网络；利用局部运算和共享权重处理空间模式。 & \hyperref[intro:19-cnn]{\ref*{intro:19-cnn}/\pageref*{intro:19-cnn}}\\
Cosine similarity & 余弦相似度 & 点积除以两向量长度之积，关注方向。 & \hyperref[intro:02-dot-product]{\ref*{intro:02-dot-product}/\pageref*{intro:02-dot-product}}\\
Cross-attention & 交叉注意力 & Query 来自一个序列，Key/Value 来自另一序列的注意力。 & \hyperref[concept:cross-attention]{\ref*{concept:cross-attention}/\pageref*{concept:cross-attention}}\\
Cross-entropy & 交叉熵 & 衡量预测概率分布与目标之间差距的常用损失。 & \hyperref[concept:cross-entropy]{\ref*{concept:cross-entropy}/\pageref*{concept:cross-entropy}}\\
Data augmentation & 数据增强 & 在标签语义不变的前提下变换样本，增加训练多样性。 & \hyperref[intro:18-generalization]{\ref*{intro:18-generalization}/\pageref*{intro:18-generalization}}\\
Decoder & 解码器 & 按任务条件产生输出序列的网络部分。 & \hyperref[intro:22-transformer-architecture]{\ref*{intro:22-transformer-architecture}/\pageref*{intro:22-transformer-architecture}}\\
Delta rule & Delta 更新规则 & 先读旧状态，用目标与旧输出的残差修正写入。 & \hyperref[concept:delta-rule]{\ref*{concept:delta-rule}/\pageref*{concept:delta-rule}}\\
Derivative & 导数 & 衡量函数在某一点的局部变化率。 & \hyperref[intro:16-loss-gradient]{\ref*{intro:16-loss-gradient}/\pageref*{intro:16-loss-gradient}}\newline \hyperref[experiment:beginner-iteration]{\ref*{experiment:beginner-iteration}/\pageref*{experiment:beginner-iteration}}\newline \hyperref[experiment:learning-rate-stability]{\ref*{experiment:learning-rate-stability}/\pageref*{experiment:learning-rate-stability}}\\
Dimension & 维度 & 向量有几个数，或矩阵有几行几列。 & \hyperref[intro:01-vectors]{\ref*{intro:01-vectors}/\pageref*{intro:01-vectors}}\\
Distillation & 蒸馏 & 把教师模型的输出或分布作为学生模型的训练信号。 & \hyperref[concept:distillation]{\ref*{concept:distillation}/\pageref*{concept:distillation}}\\
Dot product & 点积 & 对应分量相乘后求和；若向量归一化，可反映夹角余弦。 & \hyperref[intro:02-dot-product]{\ref*{intro:02-dot-product}/\pageref*{intro:02-dot-product}}\\
Dropout & 随机失活 & 训练时随机屏蔽一部分激活；评估模式下关闭随机屏蔽。 & \hyperref[concept:dropout]{\ref*{concept:dropout}/\pageref*{concept:dropout}}\\
\addlinespace[4pt]
\multicolumn{4}{@{}l}{\sffamily\bfseries\color{blue} English E--H / 英文首字母索引}\\*
Early stopping & 提前停止 & 根据验证表现决定停止训练的时机。 & \hyperref[intro:18-generalization]{\ref*{intro:18-generalization}/\pageref*{intro:18-generalization}}\\
Embedding & 嵌入 & 把离散 Token ID 映射到可学习的稠密向量。 & \hyperref[intro:20-embeddings]{\ref*{intro:20-embeddings}/\pageref*{intro:20-embeddings}}\newline \hyperref[intro:20b-task-contract]{\ref*{intro:20b-task-contract}/\pageref*{intro:20b-task-contract}}\newline \hyperref[intro:23a-from-toy-to-gpt]{\ref*{intro:23a-from-toy-to-gpt}/\pageref*{intro:23a-from-toy-to-gpt}}\\
Embedding matrix & 嵌入矩阵 & 存放词表各 Token 向量的一张权重表。 & \hyperref[intro:20-embeddings]{\ref*{intro:20-embeddings}/\pageref*{intro:20-embeddings}}\newline \hyperref[intro:20b-task-contract]{\ref*{intro:20b-task-contract}/\pageref*{intro:20b-task-contract}}\newline \hyperref[intro:23a-from-toy-to-gpt]{\ref*{intro:23a-from-toy-to-gpt}/\pageref*{intro:23a-from-toy-to-gpt}}\\
Encoder & 编码器 & 用于将输入转成上下文化表示的网络部分。 & \hyperref[intro:22-transformer-architecture]{\ref*{intro:22-transformer-architecture}/\pageref*{intro:22-transformer-architecture}}\\
Epoch & 训练轮次 & 训练过程完整遍历一次训练数据。 & \hyperref[intro:18c-training-engineering]{\ref*{intro:18c-training-engineering}/\pageref*{intro:18c-training-engineering}}\\
Evaluation contamination & 评测污染 & 测试题、答案或近重复版本进入训练数据，导致评测高估泛化。 & \hyperref[concept:evaluation-contamination]{\ref*{concept:evaluation-contamination}/\pageref*{concept:evaluation-contamination}}\newline \hyperref[concept:evaluation-contamination-benchmark]{\ref*{concept:evaluation-contamination-benchmark}/\pageref*{concept:evaluation-contamination-benchmark}}\\
Expert router & 专家路由器 & 决定当前 Token 由哪些专家处理。 & \hyperref[concept:router]{\ref*{concept:router}/\pageref*{concept:router}}\\
Exploding gradient & 梯度爆炸 & 反向传播中梯度幅度变得过大。 & \hyperref[intro:18c-training-engineering]{\ref*{intro:18c-training-engineering}/\pageref*{intro:18c-training-engineering}}\\
Feature map & 特征图 & 卷积层输出的空间数字表。 & \hyperref[intro:19-cnn]{\ref*{intro:19-cnn}/\pageref*{intro:19-cnn}}\\
Few-shot prompting & 少样本提示 & 在 Prompt 中放入少量示例以引导任务格式，是上下文学习的常见形式。 & \hyperref[concept:few-shot]{\ref*{concept:few-shot}/\pageref*{concept:few-shot}}\\
Fine-tuning & 微调 & 在已有模型权重上针对任务或行为继续训练。 & \hyperref[intro:23c-post-training]{\ref*{intro:23c-post-training}/\pageref*{intro:23c-post-training}}\\
Forward propagation & 前向传播 & 从输入沿计算图得到输出与损失。 & \hyperref[intro:16-loss-gradient]{\ref*{intro:16-loss-gradient}/\pageref*{intro:16-loss-gradient}}\newline \hyperref[experiment:beginner-iteration]{\ref*{experiment:beginner-iteration}/\pageref*{experiment:beginner-iteration}}\newline \hyperref[experiment:learning-rate-stability]{\ref*{experiment:learning-rate-stability}/\pageref*{experiment:learning-rate-stability}}\\
Fully connected layer & 全连接层 & 本层各输出单元能使用前一层所有输入特征。 & \hyperref[intro:19-cnn]{\ref*{intro:19-cnn}/\pageref*{intro:19-cnn}}\\
Function & 函数 & 输入经过确定的计算规则得到输出，例如 $y=2x+1$。 & \hyperref[intro:00-functions]{\ref*{intro:00-functions}/\pageref*{intro:00-functions}}\\
Generalization & 泛化 & 从已见样本学到能应用于新样本的规律。 & \hyperref[intro:18-generalization]{\ref*{intro:18-generalization}/\pageref*{intro:18-generalization}}\\
Gradient & 梯度 & 将各个参数方向的偏导数组合起来。 & \hyperref[intro:16-loss-gradient]{\ref*{intro:16-loss-gradient}/\pageref*{intro:16-loss-gradient}}\newline \hyperref[experiment:beginner-iteration]{\ref*{experiment:beginner-iteration}/\pageref*{experiment:beginner-iteration}}\newline \hyperref[experiment:learning-rate-stability]{\ref*{experiment:learning-rate-stability}/\pageref*{experiment:learning-rate-stability}}\\
Gradient clipping & 梯度裁剪 & 限制梯度幅度，减少过大更新的风险。 & \hyperref[intro:18c-training-engineering]{\ref*{intro:18c-training-engineering}/\pageref*{intro:18c-training-engineering}}\\
Gradient descent & 梯度下降 & 沿梯度反方向逐步调整参数，尝试降低损失。 & \hyperref[intro:16a-gradient-descent]{\ref*{intro:16a-gradient-descent}/\pageref*{intro:16a-gradient-descent}}\newline \hyperref[experiment:beginner-iteration]{\ref*{experiment:beginner-iteration}/\pageref*{experiment:beginner-iteration}}\newline \hyperref[experiment:learning-rate-stability]{\ref*{experiment:learning-rate-stability}/\pageref*{experiment:learning-rate-stability}}\\
Group Relative Policy Optimization (GRPO) & 组相对策略优化 & 以组内相对奖励构造策略更新信号；还涉及归一化和裁剪等训练步骤。 & \hyperref[concept:grpo]{\ref*{concept:grpo}/\pageref*{concept:grpo}}\\
Grouped-query attention (GQA) & 分组查询注意力 & 让一组 Query 头共享 K/V 头，降低缓存规模。 & \hyperref[concept:gqa]{\ref*{concept:gqa}/\pageref*{concept:gqa}}\\
Hallucination & 幻觉 & 模型生成流畅但事实不可靠或没有证据的内容。 & \hyperref[intro:23d-agent-systems]{\ref*{intro:23d-agent-systems}/\pageref*{intro:23d-agent-systems}}\\
Hallucination and calibration & 幻觉与校准 & 区分无依据的主张与置信分数的正确频率；在保留集逐项核验。 & \hyperref[concept:hallucination-calibration]{\ref*{concept:hallucination-calibration}/\pageref*{concept:hallucination-calibration}}\\
Hidden state & 隐藏状态 & 序列模型随位置传递的中间向量。 & \hyperref[intro:21-rnn]{\ref*{intro:21-rnn}/\pageref*{intro:21-rnn}}\\
Hyperparameter & 超参数 & 控制训练或结构的设置，如学习率、层数与批量大小。 & \hyperref[intro:18-generalization]{\ref*{intro:18-generalization}/\pageref*{intro:18-generalization}}\\
\addlinespace[4pt]
\multicolumn{4}{@{}l}{\sffamily\bfseries\color{blue} English I--L / 英文首字母索引}\\*
In-context learning (ICL) & 上下文学习 & 利用当前请求中的任务说明和示例适配输出，不直接更新模型权重。 & \hyperref[concept:in-context-learning]{\ref*{concept:in-context-learning}/\pageref*{concept:in-context-learning}}\\
Inference & 推理 & 用已训练的模型对新输入产生输出。 & \hyperref[intro:00-functions]{\ref*{intro:00-functions}/\pageref*{intro:00-functions}}\\
Key & 键向量 & 与 Query 做匹配的向量。 & \hyperref[concept:qkv]{\ref*{concept:qkv}/\pageref*{concept:qkv}}\newline \hyperref[experiment:attention-selection]{\ref*{experiment:attention-selection}/\pageref*{experiment:attention-selection}}\\
KV cache & 键值缓存 & 保存过去 Token 的 Key/Value 表示供自回归解码复用。 & \hyperref[concept:kv-cache]{\ref*{concept:kv-cache}/\pageref*{concept:kv-cache}}\newline \hyperref[concept:prefill-decode]{\ref*{concept:prefill-decode}/\pageref*{concept:prefill-decode}}\newline \hyperref[concept:gqa]{\ref*{concept:gqa}/\pageref*{concept:gqa}}\\
Label & 标签 & 监督学习中为样本提供的目标或正确类别。 & \hyperref[intro:18-generalization]{\ref*{intro:18-generalization}/\pageref*{intro:18-generalization}}\\
Learning rate & 学习率 & 控制每次梯度更新步长的系数。 & \hyperref[intro:16-loss-gradient]{\ref*{intro:16-loss-gradient}/\pageref*{intro:16-loss-gradient}}\newline \hyperref[experiment:beginner-iteration]{\ref*{experiment:beginner-iteration}/\pageref*{experiment:beginner-iteration}}\newline \hyperref[experiment:learning-rate-stability]{\ref*{experiment:learning-rate-stability}/\pageref*{experiment:learning-rate-stability}}\\
Learning rate warmup & 学习率升温 & 训练前几步由小到大调度学习率，控制初期更新步长。 & \hyperref[concept:lr-warmup]{\ref*{concept:lr-warmup}/\pageref*{concept:lr-warmup}}\\
Linear attention & 线性注意力 & 通过不同运算结构改善长序列计算或状态扩展性，读取规则由具体的核或状态更新公式定义。 & \hyperref[concept:linear-attention]{\ref*{concept:linear-attention}/\pageref*{concept:linear-attention}}\\
Linear transformation & 线性变换 & 可写成 $W\vect x$ 的运算；满足叠加规则。 & \hyperref[intro:03-matrix-multiplication]{\ref*{intro:03-matrix-multiplication}/\pageref*{intro:03-matrix-multiplication}}\\
LLM-as-a-judge & 模型评审 & 用 LLM 按评分规则评价回答，并用换序与人工抽查复核裁判。 & \hyperref[concept:llm-judge]{\ref*{concept:llm-judge}/\pageref*{concept:llm-judge}}\\
Logits & 未归一化分数 & 进入 Softmax 前、还没归一化的分数。 & \hyperref[concept:logits]{\ref*{concept:logits}/\pageref*{concept:logits}}\\
Long short-term memory & 长短期记忆网络 & 利用门控状态改进长期信息传播的一类循环网络。 & \hyperref[intro:21-rnn]{\ref*{intro:21-rnn}/\pageref*{intro:21-rnn}}\\
Long-range dependency & 长期依赖 & 较早位置的信息影响较晚位置的预测。 & \hyperref[intro:21-rnn]{\ref*{intro:21-rnn}/\pageref*{intro:21-rnn}}\\
Loss function & 损失函数 & 用数字评估预测与目标之间的差距。 & \hyperref[intro:16-loss-gradient]{\ref*{intro:16-loss-gradient}/\pageref*{intro:16-loss-gradient}}\newline \hyperref[experiment:beginner-iteration]{\ref*{experiment:beginner-iteration}/\pageref*{experiment:beginner-iteration}}\newline \hyperref[experiment:learning-rate-stability]{\ref*{experiment:learning-rate-stability}/\pageref*{experiment:learning-rate-stability}}\\
Loss mask & 损失掩码 & 选择哪些位置参与损失统计，区别于因果注意力掩码。 & \hyperref[concept:loss-mask]{\ref*{concept:loss-mask}/\pageref*{concept:loss-mask}}\\
Low-rank adaptation & 低秩适配 & 冻结多数基础权重并训练低秩参数的微调方法。 & \hyperref[concept:lora]{\ref*{concept:lora}/\pageref*{concept:lora}}\\
\addlinespace[4pt]
\multicolumn{4}{@{}l}{\sffamily\bfseries\color{blue} English M--P / 英文首字母索引}\\*
Matrix & 矩阵 & 按行和列排列的数字表；可同时执行多个点积。 & \hyperref[intro:03-matrix-multiplication]{\ref*{intro:03-matrix-multiplication}/\pageref*{intro:03-matrix-multiplication}}\\
Mean squared error & 均方误差 & 误差平方再取平均，常用于回归。 & \hyperref[intro:16-loss-gradient]{\ref*{intro:16-loss-gradient}/\pageref*{intro:16-loss-gradient}}\newline \hyperref[experiment:beginner-iteration]{\ref*{experiment:beginner-iteration}/\pageref*{experiment:beginner-iteration}}\newline \hyperref[experiment:learning-rate-stability]{\ref*{experiment:learning-rate-stability}/\pageref*{experiment:learning-rate-stability}}\\
Mini-batch & 小批次 & 将数据分为较小的批次计算梯度。 & \hyperref[intro:18c-training-engineering]{\ref*{intro:18c-training-engineering}/\pageref*{intro:18c-training-engineering}}\\
Mixed precision & 混合精度 & 不同计算或状态采用不同浮点精度，以平衡数值稳定性、存储与速度。 & \hyperref[concept:mixed-precision]{\ref*{concept:mixed-precision}/\pageref*{concept:mixed-precision}}\\
Mixture of experts (MoE) & 混合专家 & 让不同 Token 按路由选择少数专家网络进行计算。 & \hyperref[concept:moe]{\ref*{concept:moe}/\pageref*{concept:moe}}\\
Model & 模型 & 为完成任务选定的一套计算结构和参数。 & \hyperref[intro:00-functions]{\ref*{intro:00-functions}/\pageref*{intro:00-functions}}\\
Model Context Protocol & 模型上下文协议 & 模型应用连接外部工具和数据源的一种标准协议。 & \hyperref[concept:mcp]{\ref*{concept:mcp}/\pageref*{concept:mcp}}\\
Multi-head attention & 多头注意力 & 用多组可训练投影并行学习不同注意力模式。 & \hyperref[concept:multihead]{\ref*{concept:multihead}/\pageref*{concept:multihead}}\\
Multi-head latent attention (MLA) & 多头潜在注意力 & 通过潜在表示等机制尝试压缩注意力缓存。 & \hyperref[concept:mla]{\ref*{concept:mla}/\pageref*{concept:mla}}\\
Multi-layer perceptron & 多层感知机 & 多层感知机；组合矩阵乘法、偏置和非线性激活。 & \hyperref[intro:10-mlp]{\ref*{intro:10-mlp}/\pageref*{intro:10-mlp}}\\
Multi-query attention (MQA) & 多查询注意力 & 让多个 Query 头共享一组 K/V 头。 & \hyperref[concept:mqa]{\ref*{concept:mqa}/\pageref*{concept:mqa}}\\
Nonlinearity & 非线性 & 不满足简单线性叠加的变化，例如 ReLU。 & \hyperref[intro:07-functions-relu]{\ref*{intro:07-functions-relu}/\pageref*{intro:07-functions-relu}}\\
Norm & 范数 & 衡量向量长度或大小的一类函数。 & \hyperref[intro:01-vectors]{\ref*{intro:01-vectors}/\pageref*{intro:01-vectors}}\\
Normalization & 归一化 & 调整激活尺度或统计分布以辅助稳定训练。 & \hyperref[intro:22-transformer-architecture]{\ref*{intro:22-transformer-architecture}/\pageref*{intro:22-transformer-architecture}}\\
One-hot encoding & 独热编码 & 只有对应类别位置为一，其余为零的编码。 & \hyperref[intro:20-embeddings]{\ref*{intro:20-embeddings}/\pageref*{intro:20-embeddings}}\newline \hyperref[intro:20b-task-contract]{\ref*{intro:20b-task-contract}/\pageref*{intro:20b-task-contract}}\newline \hyperref[intro:23a-from-toy-to-gpt]{\ref*{intro:23a-from-toy-to-gpt}/\pageref*{intro:23a-from-toy-to-gpt}}\\
Open weights & 开放权重 & 允许获取模型权重；不必然意味着训练代码和数据完全开放。 & \hyperref[intro:23b-deployment-and-ecosystem]{\ref*{intro:23b-deployment-and-ecosystem}/\pageref*{intro:23b-deployment-and-ecosystem}}\\
Optimizer & 优化器 & 根据梯度与状态计算参数更新的算法。 & \hyperref[intro:18c-training-engineering]{\ref*{intro:18c-training-engineering}/\pageref*{intro:18c-training-engineering}}\\
Outer product & 外积 & 将列形式 Key 与行形式 Value 相乘，写成关联矩阵。 & \hyperref[concept:outer-product]{\ref*{concept:outer-product}/\pageref*{concept:outer-product}}\\
Output head & 输出头 & 把每个 Token 的隐藏向量映射为词表 logits，不等同于内部 FFN。 & \hyperref[concept:output-head]{\ref*{concept:output-head}/\pageref*{concept:output-head}}\\
Overfitting & 过拟合 & 训练表现优异但未见数据表现较差。 & \hyperref[intro:18-generalization]{\ref*{intro:18-generalization}/\pageref*{intro:18-generalization}}\\
Padding & 填充 & 在输入边界增添数值，用于调节输出大小与边界处理。 & \hyperref[intro:19-cnn]{\ref*{intro:19-cnn}/\pageref*{intro:19-cnn}}\\
Parameter & 参数 & 训练中直接学习的数，如权重 $w$ 和偏置 $b$。 & \hyperref[intro:00-functions]{\ref*{intro:00-functions}/\pageref*{intro:00-functions}}\\
Partial derivative & 偏导数 & 多变量函数对某个变量求变化率，其他变量暂时固定。 & \hyperref[intro:16-loss-gradient]{\ref*{intro:16-loss-gradient}/\pageref*{intro:16-loss-gradient}}\newline \hyperref[experiment:beginner-iteration]{\ref*{experiment:beginner-iteration}/\pageref*{experiment:beginner-iteration}}\newline \hyperref[experiment:learning-rate-stability]{\ref*{experiment:learning-rate-stability}/\pageref*{experiment:learning-rate-stability}}\\
Perplexity & 困惑度 & 平均自然对数交叉熵的指数；比较时须保持分词与计分数据口径一致。 & \hyperref[concept:perplexity]{\ref*{concept:perplexity}/\pageref*{concept:perplexity}}\newline \hyperref[concept:llm-judge]{\ref*{concept:llm-judge}/\pageref*{concept:llm-judge}}\\
Pooling & 池化 & 对局部区域做最大值或平均等汇总。 & \hyperref[intro:19-cnn]{\ref*{intro:19-cnn}/\pageref*{intro:19-cnn}}\\
Pre-Norm & 前置归一化 & 在某个计算子层之前进行归一化。 & \hyperref[concept:pre-norm]{\ref*{concept:pre-norm}/\pageref*{concept:pre-norm}}\\
Prefill and decode & 预填充与解码 & 先处理已知前缀建立缓存，再逐步处理新增 Token 并读取历史。 & \hyperref[concept:prefill-decode]{\ref*{concept:prefill-decode}/\pageref*{concept:prefill-decode}}\\
Pretraining & 预训练 & 通过大规模数据先学习一般的表示和模式。 & \hyperref[concept:pretraining]{\ref*{concept:pretraining}/\pageref*{concept:pretraining}}\\
Prompt & 提示输入 & 提供给模型的任务说明、示例、问题等提示输入。 & \hyperref[concept:prompt]{\ref*{concept:prompt}/\pageref*{concept:prompt}}\\
Prompt injection & 提示注入 & 外部内容把数据伪装成操作指令，诱导模型或应用越过原有任务与权限。 & \hyperref[concept:prompt-injection]{\ref*{concept:prompt-injection}/\pageref*{concept:prompt-injection}}\newline \hyperref[sec:agent-safety]{\ref*{sec:agent-safety}/\pageref*{sec:agent-safety}}\\
Pruning & 剪枝 & 移除或稀疏化部分权重、连接或结构，再评测质量与成本。 & \hyperref[concept:pruning]{\ref*{concept:pruning}/\pageref*{concept:pruning}}\\
\addlinespace[4pt]
\multicolumn{4}{@{}l}{\sffamily\bfseries\color{blue} English Q--T / 英文首字母索引}\\*
Quantization & 量化 & 用更低数值精度保存或计算部分权重与激活。 & \hyperref[concept:quant]{\ref*{concept:quant}/\pageref*{concept:quant}}\\
Query & 查询向量 & 查询向量，表达当前这个位置想寻找什么。 & \hyperref[concept:qkv]{\ref*{concept:qkv}/\pageref*{concept:qkv}}\newline \hyperref[experiment:attention-selection]{\ref*{experiment:attention-selection}/\pageref*{experiment:attention-selection}}\\
Recurrent neural network & 循环神经网络 & 循环神经网络；用递推隐藏状态处理序列。 & \hyperref[intro:21-rnn]{\ref*{intro:21-rnn}/\pageref*{intro:21-rnn}}\\
Regularization & 正则化 & 通过参数惩罚或其他限制降低过拟合风险。 & \hyperref[intro:18-generalization]{\ref*{intro:18-generalization}/\pageref*{intro:18-generalization}}\\
Reinforcement learning from human feedback & 人类反馈强化学习 & 用人类反馈相关的偏好信号辅助模型行为优化。 & \hyperref[concept:rlhf]{\ref*{concept:rlhf}/\pageref*{concept:rlhf}}\\
Reinforcement learning with verifiable rewards (RLVR) & 可验证奖励强化学习 & 利用可程序化核验的任务奖励进行学习。 & \hyperref[concept:rlvr]{\ref*{concept:rlvr}/\pageref*{concept:rlvr}}\\
Residual connection & 残差连接 & 把输入直接加回模块输出，为信息传播提供旁路。 & \hyperref[concept:residual]{\ref*{concept:residual}/\pageref*{concept:residual}}\\
Retrieval-augmented generation & 检索增强生成 & 检索外部资料并将结果提供给生成模型。 & \hyperref[concept:rag]{\ref*{concept:rag}/\pageref*{concept:rag}}\\
Reusable skill & 可复用技能 & 把一类任务的说明、约束和脚本包装为可复用能力。 & \hyperref[concept:skill]{\ref*{concept:skill}/\pageref*{concept:skill}}\\
Reward model & 奖励模型 & 从评价或偏好数据学习对候选回答的评分，不保证评分等同事实真值。 & \hyperref[concept:reward-model]{\ref*{concept:reward-model}/\pageref*{concept:reward-model}}\\
Root Mean Square Normalization (RMSNorm) & 均方根归一化 & 按特征的均方根控制数值尺度。 & \hyperref[concept:rmsnorm]{\ref*{concept:rmsnorm}/\pageref*{concept:rmsnorm}}\\
Rotary position embedding (RoPE) & 旋转位置编码 & 对 Q/K 中的分量对做位置相关旋转。 & \hyperref[concept:rope]{\ref*{concept:rope}/\pageref*{concept:rope}}\\
Sample & 样本 & 用于训练或评估的一组输入及可能的目标值。 & \hyperref[intro:00-functions]{\ref*{intro:00-functions}/\pageref*{intro:00-functions}}\\
Scaled dot-product attention & 缩放点积 & 把 Query 与 Key 的点积按维度尺度缩放。 & \hyperref[concept:scaled-dot-product]{\ref*{concept:scaled-dot-product}/\pageref*{concept:scaled-dot-product}}\\
Self-attention & 自注意力 & Query、Key、Value 都来源于同一序列的注意力机制。 & \hyperref[concept:self-attention]{\ref*{concept:self-attention}/\pageref*{concept:self-attention}}\\
Sequence & 序列 & 有顺序的输入元素集合，例如一串 Token。 & \hyperref[intro:20-embeddings]{\ref*{intro:20-embeddings}/\pageref*{intro:20-embeddings}}\newline \hyperref[intro:20b-task-contract]{\ref*{intro:20b-task-contract}/\pageref*{intro:20b-task-contract}}\newline \hyperref[intro:23a-from-toy-to-gpt]{\ref*{intro:23a-from-toy-to-gpt}/\pageref*{intro:23a-from-toy-to-gpt}}\\
Shifted targets & 目标位置错位 & 每个输入位置的目标为下一位置 Token。 & \hyperref[concept:shifted-targets]{\ref*{concept:shifted-targets}/\pageref*{concept:shifted-targets}}\\
Sliding-window attention & 滑动窗口注意力 & 每个位置只读取一个有限的可见前缀窗口。 & \hyperref[concept:sliding-window-attention]{\ref*{concept:sliding-window-attention}/\pageref*{concept:sliding-window-attention}}\\
Softmax & Softmax 归一化 & 将分数转成非负且总和为一的一组权重。 & \hyperref[concept:softmax]{\ref*{concept:softmax}/\pageref*{concept:softmax}}\\
Sparse attention & 稀疏注意力 & 只计算选择的 Token 对，不要求全部两两比较。 & \hyperref[concept:sparse-attention]{\ref*{concept:sparse-attention}/\pageref*{concept:sparse-attention}}\\
Stride & 步幅 & 卷积核在相邻位置移动的距离。 & \hyperref[intro:19-cnn]{\ref*{intro:19-cnn}/\pageref*{intro:19-cnn}}\\
Subagent and agent-to-agent (A2A) & 子智能体与 A2A 通信 & 前者为子任务分配独立执行单元；后者规范独立 Agent 间通信。 & \hyperref[concept:subagent]{\ref*{concept:subagent}/\pageref*{concept:subagent}}\\
Supervised fine-tuning & 监督微调 & 利用示范输入输出对继续调整模型。 & \hyperref[concept:sft]{\ref*{concept:sft}/\pageref*{concept:sft}}\\
SwiGLU & SwiGLU 门控 & 使用两个投影分支与逐元素门控的 FFN 变体。 & \hyperref[concept:swiglu]{\ref*{concept:swiglu}/\pageref*{concept:swiglu}}\\
Task contract & 任务契约 & 事先约定输入、目标、张量尺寸和结果的含义。 & \hyperref[concept:task-contract]{\ref*{concept:task-contract}/\pageref*{concept:task-contract}}\\
Teacher forcing & 教师强制 & 训练时以已知真实序列作为条件，一次预测多个后续位置。 & \hyperref[concept:teacher-forcing]{\ref*{concept:teacher-forcing}/\pageref*{concept:teacher-forcing}}\\
Temperature & 温度 & 通过缩放 logits 控制输出分布尖锐程度的参数。 & \hyperref[concept:sampling]{\ref*{concept:sampling}/\pageref*{concept:sampling}}\\
Test set & 测试集 & 尽量保留到最后，估计泛化性能的数据。 & \hyperref[intro:18-generalization]{\ref*{intro:18-generalization}/\pageref*{intro:18-generalization}}\\
Time step & 时间步 & RNN 递推中当前处理到的序列位置。 & \hyperref[intro:21-rnn]{\ref*{intro:21-rnn}/\pageref*{intro:21-rnn}}\\
Token & 词元 & 模型分词器定义的文本处理单元，不一定是一个完整词。 & \hyperref[intro:20a-tokenizer]{\ref*{intro:20a-tokenizer}/\pageref*{intro:20a-tokenizer}}\newline \hyperref[experiment:real-chinese-tokenizer]{\ref*{experiment:real-chinese-tokenizer}/\pageref*{experiment:real-chinese-tokenizer}}\\
Tokenizer & 分词器 & 将文本切成 Token 并映射到整数 ID 的方法。 & \hyperref[intro:20a-tokenizer]{\ref*{intro:20a-tokenizer}/\pageref*{intro:20a-tokenizer}}\newline \hyperref[experiment:real-chinese-tokenizer]{\ref*{experiment:real-chinese-tokenizer}/\pageref*{experiment:real-chinese-tokenizer}}\\
Tool calling & 工具调用 & 模型提出结构化请求，系统在检查后调用外部工具。 & \hyperref[concept:tool-calling]{\ref*{concept:tool-calling}/\pageref*{concept:tool-calling}}\newline \hyperref[elective:real-function-call]{\ref*{elective:real-function-call}/\pageref*{elective:real-function-call}}\newline \hyperref[sec:agent-safety]{\ref*{sec:agent-safety}/\pageref*{sec:agent-safety}}\\
Top-k sampling & Top-k 采样 & 只在概率最高的若干候选中继续采样。 & \hyperref[concept:sampling]{\ref*{concept:sampling}/\pageref*{concept:sampling}}\\
Top-p sampling & Top-p 采样 & 保留累积概率达到阈值的一组高概率候选。 & \hyperref[concept:sampling]{\ref*{concept:sampling}/\pageref*{concept:sampling}}\\
Training & 训练 & 利用数据和损失函数更新模型参数。 & \hyperref[intro:00-functions]{\ref*{intro:00-functions}/\pageref*{intro:00-functions}}\\
Training set & 训练集 & 用于直接计算参数梯度与更新权重的数据。 & \hyperref[intro:18-generalization]{\ref*{intro:18-generalization}/\pageref*{intro:18-generalization}}\\
Transformer & Transformer 架构 & 把注意力、逐位置 MLP、残差和归一化等组合的架构。 & \hyperref[intro:22-transformer-architecture]{\ref*{intro:22-transformer-architecture}/\pageref*{intro:22-transformer-architecture}}\\
Transpose & 转置 & 交换行与列，例如 $V$ 转成 $V^\top$。 & \hyperref[intro:03-matrix-multiplication]{\ref*{intro:03-matrix-multiplication}/\pageref*{intro:03-matrix-multiplication}}\\
\addlinespace[4pt]
\multicolumn{4}{@{}l}{\sffamily\bfseries\color{blue} English U--Z / 英文首字母索引}\\*
Underfitting & 欠拟合 & 模型没学到主要规律，训练和评估表现都不理想。 & \hyperref[intro:18-generalization]{\ref*{intro:18-generalization}/\pageref*{intro:18-generalization}}\\
Validation set & 验证集 & 用于选择超参数、模型方案和提前停止的数据。 & \hyperref[intro:18-generalization]{\ref*{intro:18-generalization}/\pageref*{intro:18-generalization}}\\
Value & 值向量 & 被注意力权重组合后形成输出的向量。 & \hyperref[concept:qkv]{\ref*{concept:qkv}/\pageref*{concept:qkv}}\newline \hyperref[experiment:attention-selection]{\ref*{experiment:attention-selection}/\pageref*{experiment:attention-selection}}\\
Vanishing gradient & 梯度消失 & 反向传播到较早层时梯度变得过小。 & \hyperref[intro:18c-training-engineering]{\ref*{intro:18c-training-engineering}/\pageref*{intro:18c-training-engineering}}\\
Vector & 向量 & 有顺序的一列数字；可用来表示一个对象的多个特征。 & \hyperref[intro:01-vectors]{\ref*{intro:01-vectors}/\pageref*{intro:01-vectors}}\\
Vector retrieval & 向量检索 & 通过查询向量和候选向量的相似性寻找内容。 & \hyperref[intro:23d-agent-systems]{\ref*{intro:23d-agent-systems}/\pageref*{intro:23d-agent-systems}}\\
Vision-language alignment & 图文对齐 & 通过配对或回答训练，让视觉和文字表示对应同一内容。 & \hyperref[concept:vision-language-alignment]{\ref*{concept:vision-language-alignment}/\pageref*{concept:vision-language-alignment}}\\
Vocabulary & 词表 & 允许的 Token 与其 ID 的集合。 & \hyperref[intro:20a-tokenizer]{\ref*{intro:20a-tokenizer}/\pageref*{intro:20a-tokenizer}}\newline \hyperref[experiment:real-chinese-tokenizer]{\ref*{experiment:real-chinese-tokenizer}/\pageref*{experiment:real-chinese-tokenizer}}\\
Weight & 权重 & 控制输入特征对输出贡献的参数。 & \hyperref[intro:00-functions]{\ref*{intro:00-functions}/\pageref*{intro:00-functions}}\\
Workflow & 工作流 & 把多个固定或动态步骤组合成一个业务过程。 & \hyperref[concept:workflow]{\ref*{concept:workflow}/\pageref*{concept:workflow}}\\
\end{longtable}
\endgroup

\subsubsection*{不要混淆的五对概念}
\begin{itemize}
\item \textbf{参数与超参数}：前者主要由训练算法学习，后者控制模型或训练设计。
\item \textbf{训练与推理}：训练会更新参数，普通推理使用固定参数。
\item \textbf{模型与 Agent}：模型负责计算预测，Agent 系统还涉及工具、状态与动作。
\item \textbf{Embedding 与检索}：向量表示是检索的一种工具，检索本身还需索引、过滤和验证。
\item \textbf{开放权重与开源}：开放模型参数不等于所有训练数据、代码和许可均公开。
\end{itemize}
